% -*- TeX-master: "../main.tex" -*-

\chapter{Validación}\label{sec:validacion}



El fin de este capítulo es presentar un conjunto de casos que cumplen
un doble propósito: {(1)} aportar evidencia de que el intérprete
implementa correctamente el lenguaje GPLC y {(2)}  ilustrar el
uso práctico del lenguaje mediante programas pequeños pero
representativos.  La evidencia obtenida se basa en el contraste entre
los resultados producidos por el intérprete y el comportamiento
esperado: en los casos exitosos, las distribuciones retornadas y sus
valores con ascripción, y en los casos fallidos, los diagnósticos de
error reportados. Como complemento, al final del capítulo se incluye
una subsección de limitaciones que sintetiza dificultades observadas
durante la implementación y que motivan trabajo futuro.

A diferencia de una verificación formal completa, la validación que se
presenta aquí es funcional y dirigida por casos.  Cada programa está
diseñado para cubrir una expresión específica (o una interacción
acotada entre expresiones) y contrastar el comportamiento observado
con las reglas y decisiones de diseño descritas en capítulos previos.
Este enfoque resulta adecuado para GPLC, ya que lo que se busca
contrastar se refleja en hechos directamente observables: la
aceptación de programas bien formados y la producción de
distribuciones con valores y tipos coherentes; la propagación de
probabilidades concretas y desconocidas de acuerdo con las reglas del
lenguaje; y el rechazo estático o dinámico de programas inválidos
mediante diagnósticos precisos con localización y explicación de la
inconsistencia detectada.

Finalmente, estos casos también cumplen un rol ilustrativo, ya que
muestran patrones de uso propuestos, ejemplos mínimos y casos de borde
que aclaran qué se acepta, qué se rechaza y por qué. Por ello, el
capítulo está organizado en torno a familias de expresiones, agrupando
los programas por la construcción que se desea validar.


\section{Metodología}

Para cada caso se documentan tres elementos: {(1)} el código
fuente GPLC, {(2)} la salida del intérprete (ya sea una
distribución simbólica o un mensaje de error), y {(3)} una
explicación breve del aspecto que se está validando. En los casos
exitosos, la validación se apoya en observar que: {(i)} el
programa es aceptado por el chequeo estático, {(ii)} la
distribución retornada contiene los valores esperados, y
{(iii)} los tipos que figuran en cada valor son los
correctos. En los casos fallidos, se espera un rechazo estático o
dinámico según corresponda, y se valida que el intérprete reporte el
error con un mensaje informativo y una localización adecuada.

Los ejemplos se seleccionan además para cubrir el \emph{carácter
  gradual} del lenguaje: probabilidades desconocidas \verb|?|, el tipo
desconocido \verb|?|, y el efecto de estos elementos en consistencia,
ascripción y construcción de distribuciones. Con esto, el capítulo
busca evidenciar no solo que el intérprete evalúa programas, sino
también que preserva las restricciones semánticas y de integridad que
motivan el diseño de GPLC para escenarios con incertidumbre e
información parcial.



\section{Casos}

Los casos se organizan por construcción sintáctica: selección
probabilística para distribuir probabilidad entre dos ramas,
\lstinline|let| para definir variables, lambdas para abstracción y
aplicación, \lstinline|if| para control y ascripción para imponer
tipos. Cada conjunto de casos reúne programas mínimos que aíslan el
comportamiento bajo estudio, y luego incorpora ejemplos donde la
interacción entre componentes graduales (tipos y probabilidades
desconocidas) es relevante. Cuando corresponde, también se incluyen
casos de rechazo estático y dinámico para ilustrar las condiciones
bajo las cuales el intérprete protege la integridad semántica de
probabilidades y distribuciones.

\subsection{Selección binaria probabilística} El Código
\ref{lst:choice1} muestra un caso elemental de selección
probabilística: la selección de un número real con probabilidad
concreta 0.5. La salida muestra una distribución simbólica donde los
elementos 1 y 0 están ascritos por el tipo \lstinline|Real|, y
ponderados por probabilidades simbólicas.


\begin{gplcblock}[caption={Selección de números reales},label={lst:choice1}]
0.5 $ 1 : 0

>>>

{
	 [Real] 1. :: Real : SymProb19 ,
	 [Real] 0. :: Real : SymProb20 
}   
\end{gplcblock}




El Código \ref{lst:choice2} muestra un caso de selección con ramas
no homogéneas: valores de tipo \lstinline|Real| y \lstinline|Bool| son
seleccionados, lo cual constituye una decisión de diseño importante:
Las ramas no necesitan ser de un mismo tipo. Además la selección es
hecha con una probabilidad incógnita. La variable simbólica libre
asociada a dicha probabilidad queda internamente codificada en la
fórmula de la distribución retornada.

  
\begin{gplcblock}[caption={Selección de booleano y número},label={lst:choice2}]
? $ true : 10

>>>

{
	 [Bool] true :: Bool : SymProb19 ,
	 [Real] 10. :: Real : SymProb20 
}  
\end{gplcblock}

El Código \ref{lst:choice3} muestra un ejemplo de selección anidada,
es decir, una selección con una rama que, a su vez,
es una selección. En la salida se observa cómo los
tres valores conforman el soporte de la distribución
retornada.

\begin{gplcblock}[caption={Selección anidada},label={lst:choice3}]
0.4 $ (? $ true : false) : 10

>>>

{
	 [Bool] true :: Bool : SymProb71 ,
	 [Bool] false :: Bool : SymProb72 ,
	 [Real] 10. :: Real : SymProb73 
}  
\end{gplcblock}



En el Código \ref{lst:choice4}, la expresión de selección tiene una
ascripción de tipo de distribución que verifica estáticamente su
consistencia con el tipo inferido de la expresión. En la distribución
ascrita, el tipo \lstinline|Real| concentra el peso total de los
elementos reales, lo cual comprueba que en la relación de consistencia
la distribución de los pesos de un tipo no afecta la consistencia
entre las distribuciones. Por lo tanto, el resultado de un juicio de
consistencia está determinado estrictamente por relaciones semánticas
entre los tipos, sin importar la sintaxis específica. Esto es cierto
en general para todas las relaciones y operaciones entre
distribuciones.

\begin{gplcblock}[caption={Selección con ascripción},label={lst:choice4}]
(0.5 $ 1 : (0.4 $ 0 : false)) :: {Real : 0.7 , Bool : 0.3}

>>>

{
	 [Real] 1. :: Real : SymProb46 ,
	 [Real] 0. :: Real : SymProb47 ,
	 [Bool] false :: Bool : SymProb48 
}  
\end{gplcblock}
  
Con un carácter más gradual, los Códigos \ref{lst:choice5.1},
\ref{lst:choice5.2} y \ref{lst:choice5.3} muestran la flexibilidad que
se logra con el tipo desconocido \lstinline|?| y la probabilidad
desconocida \lstinline|?|.  El Código \ref{lst:choice5.1} muestra
que el tipo \lstinline|?| --consistente con los tipos \lstinline|Real|
y \lstinline|Bool|-- puede distribuir su peso entre \lstinline|Real| y
\lstinline|Bool| para lograr la consistencia con el tipo de la
expresión. El Código \ref{lst:choice5.2} evidencia que la
distribución con probabilidades desconocidas es consistente con toda
expresión cuyo tipo sea una distribución de \lstinline|Real| y
\lstinline|Bool|, o incluso una distribución con solo uno de esos dos
tipos, ya que la probabilidad \lstinline|?| es compatible con 0.0, es
decir, puede representar también la probabilidad de un tipo con peso
nulo. Finalmente, el Código \ref{lst:choice5.3} demuestra que se
puede anotar la selección con cualquier distribución de
\lstinline|Real| y \lstinline|Bool|, dado que la selección con
probabilidad desconocida es consistente con distribución de pesos en
tipo de ascripción.


\begin{gplcblock}[caption={Distribución con tipo no especificado},label={lst:choice5.1}]
(0.3 $ true : (-1)) :: {Bool : 0.2, ? : 0.8}

>>>

{
	 [Bool] true :: Bool : SymProb21 ,
	 [Bool] true :: ? : SymProb22 ,
	 [Real] -1. :: ? : SymProb23 
}  
\end{gplcblock}

\begin{gplcblock}[caption={Distribución con probabilidades no especificadas},label={lst:choice5.2}]
(0.3 $ true : (-1)) :: {Real : ?, Bool : ?}

>>>

{
	 [Bool] true :: Bool : SymProb21 ,
	 [Real] -1. :: Real : SymProb22 
}  
\end{gplcblock}

\begin{gplcblock}[caption={Ascripción de selección con prob. \texttt{?}},label={lst:choice5.3}]
(? $ true : (-1)) :: {Real : 0.5 , Bool : 0.5}

>>>

{
	 [Bool] true :: Bool : SymProb21 ,
	 [Real] -1. :: Real : SymProb22 
}  
\end{gplcblock}



Para demostrar con claridad el funcionamiento efectivo del mecanismo
estático de tipado, se analizan un par de casos donde la ascripción
impuesta a la selección gatilla un error estático.

En el Código \ref{lst:choice6}, la ponderación exacta que tiene el
tipo \lstinline|Bool| en la distribución inducida por la selección, es
exactamente \lstinline|0.4|. Sin embargo, en la ascripción, el tipo
`Bool` tiene una ponderación más alta.


\begin{gplcblock}[caption={Inconsistencia por prob. de \texttt{Bool}},label={lst:choice6}]
(0.4 $ false : lambda x:?.x) :: {Bool:0.45 , ?:?}

>>>

Static Error [line 1, cols 1-49]: In distribution ascription, type 0.4 * ({Bool:1}) + 0.6 * ({? -> {?:1}:1}) of term inconsistent with ascription type {Bool:0.45, ?:?}
\end{gplcblock}

Para terminar, el Código \ref{lst:choice7} falla porque en la
anotación no hay ningún tipo consistente con el tipo funcional
\lstinline|Real -> {Real:1}| que represente su peso \lstinline|0.5| en
el tipo de distribución ascrito.




\begin{gplcblock}[caption={Inconsistencia por falta de tipo funcional},label={lst:choice7}]
(0.5 $ lambda x:Real.1 : false) :: {Bool : ? , Real : ?}

>>>

Static Error [line 1, cols 1-56]: In distribution ascription, type 0.5 * ({Real -> {Real:1}:1}) + 0.5 * ({Bool:1}) of term inconsistent with ascription type {Bool:?, Real:?}
\end{gplcblock}

\subsection{Let}
La semántica de la operación \lstinline|let| consiste en asignar a la
variable \lstinline|x| cada uno de los elementos contenidos en la
distribución implícita de la expresión nombrada, y luego establecer
una distribución de probabilidad que contenga las instancias del
cuerpo con cada versión de variable \lstinline|x|. La probabilidad de
cada valor asignado a \lstinline|x|, dentro de la distribución ligada,
se toma como el peso del cuerpo en la distribución resultante.

En el primer ejemplo (Código \ref{lst:let1}), el programa es semánticamente
equivalente a la expresión \lstinline|(0.3 $ true : 1)|. La
distribución de valores retornada muestra la combinación
probabilística efectiva del cuerpo del \verb|let| (simplemente la variable
\verb|x|), evaluado en todo el rango de la variable. La variable funciona como en la mayoría de los lenguajes cuando la
expresión nombrada representa solo un elemento. Por ejemplo, en los casos
\lstinline|y = false| o \lstinline|id = lambdax:?.x|.

\begin{gplcblock}[caption={Definición de variable local},label={lst:let1}]
let x = (0.3 $ true : 1) in x

>>>

{
	 [Bool] true :: Bool : SymProb41 ,
	 [Real] 1. :: Real : SymProb42 
}  
\end{gplcblock}

El mayor potencial del bloque \lstinline|let| consiste naturalmente en
la operabilidad que le otorga a los elementos de una variable de
distribución.  En el Código \ref{lst:let2.1}, la variable b es el
argumento de una operación lógica, y en el Código \ref{lst:let2.2}
la variable r participa en una operación aritmética. Aunque no sea
visible en los pesos simbólicos, la variable propaga a la distribución
resultante el peso probabilístico de los valores individuales
asignados a ella.


\begin{gplcblock}[caption={Uso de variable en operación lógica},label={lst:let2.1}]
let b = (0.9 $ true : false) in (true or b)

>>>

{
	 [Bool] true :: Bool : SymProb73 ,
	 [Bool] true :: Bool : SymProb74 
}  
\end{gplcblock}



\begin{gplcblock}[caption={Uso de variable en operación aritmética},label={lst:let2.2}]
let r = (? $ 0 : 10) in (7 * r)

>>>

{
	 [Real] 0. :: Real : SymProb73 ,
	 [Real] 70. :: Real : SymProb74 
}  
\end{gplcblock}

El Código \ref{lst:let3} muestra bloques \lstinline|let| anidados.
Dada la naturaleza semántica de la expresión \lstinline|let|, la
expresión final en términos de \lstinline|n| y \lstinline|b| combina
todos los valores posibles de las variables, cuatro combinaciones en
total. Además, el programa se verifica con la ascripción de un tipo de
distribución consistente con la equiprobabilidad de \lstinline|Real| y
\lstinline|Bool|.




\begin{gplcblock}[caption={Bloques \texttt{let} anidados},label={lst:let3}]
(let n = (? $ 1 : 7) in
let b = (0.3 $ true : false) in
(0.5 $ 21 / n : b and false))  :: {Real:0.5, Bool:0.5}

>>>

{
	 [Real] 21. :: Real : SymProb403 ,
	 [Bool] false :: Bool : SymProb404 ,
	 [Real] 21. :: Real : SymProb405 ,
	 [Bool] false :: Bool : SymProb406 ,
	 [Real] 3. :: Real : SymProb407 ,
	 [Bool] false :: Bool : SymProb408 ,
	 [Real] 3. :: Real : SymProb409 ,
	 [Bool] false :: Bool : SymProb410 
}  
\end{gplcblock}

El ejemplo del Código~\ref{lst:let4} exhibe elementos importantes de la expresión
\lstinline|let| en conexión con las expresiones \lstinline|lambda|. Es
conveniente la asignación de una lambda a una variable, pues
permite una aplicación mucho menos verbosa y más práctica que la
aplicación de una lambda literal directamente al argumento. En segundo
lugar, el programa demuestra la capacidad del bloque \lstinline|let|
para extraer funciones de una selección.  La ascripción verifica que
los resultados son exclusivamente de tipo \lstinline|Real|.


\begin{gplcblock}[caption={Asignación de lambda a variable},label={lst:let4}]
(let step_forward = lambdax.x+1 in
let step_back = lambday.y-1 in
let f = (? $ step_forward : step_back) in
(f 0)) :: {Real:1.0}

>>>

{
	 [Real] 1. :: Real : SymProb145 ,
	 [Real] -1. :: Real : SymProb146 
}  
\end{gplcblock}

En el contexto de los bloques \lstinline|let| y las expresiones
lambda, es pertinente mencionar que si se usa una variable que no ha
sido definida con los medios que la sintaxis ofrece, como la
asignación de una expresión en un bloque \lstinline|let|, o como
parámetro formal de una expresión lambda, entonces el intérprete
gatilla un error cuyo mensaje especifica cuál es la variable no
definida, y la ubicación exacta del código en que fue utilizada.  En
el Código \ref{lst:invalid-var} se muestra el uso inválido de la
variable \lstinline|x| en el cuerpo de la expresión lambda.

\begin{gplcblock}[caption={Uso de variable no definida},label={lst:invalid-var}]
lambda y : Bool . (x + y)

>>>

Static Error [line 1, cols 19-20]: variable x not found
\end{gplcblock}


\subsection{Lambda}

\begin{comment}
Las expresiones lambda tienen la sintaxis
\lstinline|lambda<var>:<simple>.<expr>|.  La
sintaxis\lstinline| lambda<var>.<expr>| es azúcar sintático de
\lstinline|lambda<var>:?.<expr>|. Si bien el argumento de una función
debe ser una expresión de tipo simple, es decir, un valor (número,
booleano, variable, o lambda), el cuerpo de la función puede ser
cualquier clase de expresión.
\end{comment}

Las expresiones lambda tienen la sintaxis
\lstinline|lambda <var> : <simple> . <expr>|. La
sintaxis \lstinline|lambda <var> . <expr>| es
azúcar sintáctico para
\lstinline|lambda <var> : ? . <expr>|.  Si bien el argumento de una función
debe ser una expresión de tipo simple, es decir, un valor (número,
booleano, variable, o lambda), el cuerpo de la función puede ser
cualquier clase de expresión.

En el programa del Código~\ref{lst:lambda1.1}, la función
\lstinline|seven_untyped| tiene un parámetro formal de tipo
\lstinline|?|.  Por lo tanto, la función es capaz de recibir un
argumento con cualquier tipo simple, sin levantar un error estático de
inconsistencia, porque \lstinline|?| es consistente con todo tipo
simple. (En el caso particular de esta función, tampoco sería posible
tener un error dinámico, porque el cuerpo no tiene la variable como
subexpresión). Los argumentos 1 y true se pasan exitosamente y el
programa termina.

\begin{gplcblock}[caption={Función no tipada},label={lst:lambda1.1}]
let seven_untyped = lambdax:?.7 in
let a1 = (seven_untyped 1) in
let a2 = (seven_untyped true) in
a1 * a2 

>>>

{
	 [Real] 49. :: Real : SymProb85 
}  
\end{gplcblock}

En cambio, en el programa del Código~\ref{lst:lambda1.2}, la lambda
\lstinline|seven_typed| se define con un parámetro de tipo
\lstinline|Real|, explícitamente declarado. En este segundo ejemplo,
el tipado gradual rechaza estáticamente la expresión
\lstinline|(seven_typed false)|, debido a la inconsistencia del tipo
del argumento (\lstinline|Bool|) con el tipo del parámetro
(\lstinline|Real|).

Por lo tanto, el sistema gradual prueba en esta implementación tener
la flexibilidad para recibir argumentos de todo tipo, cuando el tipo
del parámetro se omite o se declara como \lstinline|?|, y decide en
tiempo de ejecución la compatibilidad del argumento. Por otro lado, el
grado de refinamiento del tipo declarado del parámetro puede resultar
en un rechazo estático del programa.




\begin{gplcblock}[caption={Función explícitamente tipada},label={lst:lambda1.2}]
let seven_typed = lambdax:Real.7 in
let b1 = (seven_typed false) in
let b2 = (seven_typed 0) in
b1 + b2  

>>>

Static Error [line 2, cols 9-28]: In function, expected argument of type Real, but found Bool
\end{gplcblock}


En el lenguaje GPLC, las funciones son de primera clase, es decir, son
valores que pueden a su vez pasarse como argumento a otra
función. En el Código~\ref{lst:lambda2}, se define un aplicador como
una función que recibe otra función como argumento, y la aplica al
número \lstinline|5|. Luego, se define \lstinline|sqr| como función a
ser aplicada. En la expresión principal, el aplicador toma
\lstinline|sqr| como función argumento y la aplica al \lstinline|5|.

\begin{gplcblock}[caption={Función como argumento de otra función},label={lst:lambda2}]
let apply = lambda f : Real -> {Real:1} . (f 5) in
let sqr = lambda r . (r * r) in
(apply sqr)

>>>

{
	 [Real] 25. :: Real : SymProb72 
}  

\end{gplcblock}





En el Código~\ref{lst:lambda3} se demuestra cómo una función puede
operar sobre una variable ligada a los elementos de una distribución
inducida por una selección probabilística.  En la evaluación del
programa, el \lstinline|64| es dividido por
\lstinline|2|,\lstinline|4|, y \lstinline|8|, distribuyendo el peso
total de una selección numérica (0.8) de acuerdo con el peso de cada valor
de \lstinline|n|.

\begin{gplcblock}[caption={Función aplicada a identificador de distribución},label={lst:lambda3}]
let n = 0.5 $ ( ? $ 2 : 4 ) : 8 in
let f = (lambda x : Real . (0.2 $ false : 64/x)) :: Real -> {Bool:0.2, Real:0.8} in
(f n)

>>>
{
	 [Bool] false :: Bool : SymProb362 ,
	 [Real] 32. :: Real : SymProb363 ,
	 [Bool] false :: Bool : SymProb364 ,
	 [Real] 16. :: Real : SymProb365 ,
	 [Bool] false :: Bool : SymProb366 ,
	 [Real] 8. :: Real : SymProb367 
}  
\end{gplcblock}



\subsection{Condicional}

El bloque condicional en este lenguaje tiene
características similares a las que presenta
en otros lenguajes: El tipo de la condición
debe tener un tipo consistente con \lstinline|Bool|, y
los tipos de las ramas deben ser iguales, a
diferencia de la heterogeneidad que pueden tener
las ramas de la selección binaria probabilística.
En el programa del Código~\ref{lst:if1}, la condición \lstinline|false| retorna \lstinline|-1|,
y ambas ramas tienen el mismo tipo.


\begin{gplcblock}[caption={Condicional},label={lst:if1}]
if false then 1 else (-1)

>>>

{
	 [Real] -1. :: Real : SymProb11 
}  
\end{gplcblock}

En cambio, el programa \ref{lst:if2} es rechazado
estáticamente debido a que las ramas de
ejecución tienen tipos diferentes.


\begin{gplcblock}[caption={Condicional fallido},label={lst:if2}]
if true then 1 else true

>>>

Static Error [line 1, cols 0-20]: In (if), branches have unequal types {Real:1} and {Bool:1}
\end{gplcblock}


Los Códigos \ref{lst:if3} y \ref{lst:if4} comprueban que en
tiempo de ejecución no se ejecutan ambas ramas.
En el programa \ref{lst:if3}, se retorna exitosamente la
rama \lstinline|true|. En la rama false, el argumento de \lstinline|f|
pasa el chequeo estático de consistencia, porque
\lstinline|Bool| es consistente con el tipo \lstinline|?| del parámetro
formal. Sin embargo, en el programa \ref{lst:if4}, la rama
false se ejecuta, y en tiempo de ejecución, el
programa es rechazado dinámicamente.

\begin{gplcblock}[caption={Evaluación de rama sin error},label={lst:if3}]
let double = lambda x : ? . (x + x) in
if true then (double 1) else (double true) 

>>>

{
	 [Real] 2. :: Real : SymProb59 
}  
\end{gplcblock}


\begin{gplcblock}[caption={Evaluación de rama con error},label={lst:if4}]
let double = lambda x : ? . (x + x) in
if false then (double 1) else (double true)

>>>

Run-Time Error [line 1, cols 25-26]: expected value Real; found Bool
\end{gplcblock}

El código \ref{lst:if5} demuestra dos aspectos relevantes del lenguaje y
su implementación. En primer lugar, la condición del bloque
\lstinline|if| puede ser una variable que recorre los elementos de una
distribución. En segundo lugar, el ejemplo muestra una forma efectiva
de implementar una selección binaria que fuerza la homogeneidad de las
ramas. En general,
\lstinline|let cond = (p $ true : false) in (if cond then m else n)| retorna una distribución con un solo tipo (distribuido en distintos valores).


\begin{gplcblock}[caption={Condición variable},label={lst:if5}]
let condition = (0.7 $ true : false) in
if condition then 100 else 200

>>>

{
	 [Real] 100. :: Real : SymProb59 ,
	 [Real] 200. :: Real : SymProb60 
}  
\end{gplcblock}


  

  
\subsection{Ascripción simple}
La operación de ascripción con un tipo simple tiene las siguientes
funciones: (1) Verificar estáticamente la consistencia del tipo
inferido de la subexpresión con el tipo ascrito, y (2) Realizar una
conversión del tipo de la subexpresión al tipo ascrito, ya que
estáticamente, el tipo de la expresión \lstinline|<expr>::<type>| es
el tipo simple ascrito.


El programa \ref{lst:ascr1} termina exitosamente y retorna una
distribución de Dirac, es decir, una distribución cuyo soporte
contiene un único valor con probabilidad 1. Por otro lado el programa
\ref{lst:ascr2} falla estáticamente, porque la expresión tiene tipo
\lstinline|Bool|, inconsistente con \lstinline|Real|. Finalmente, en
el programa \ref{lst:ascr3}, el valor \lstinline|true|, ascrito con
\lstinline|?| se asigna a la variable \lstinline|unknown|.  Esta
ascripción funciona efectivamente como una conversión de tipo, que deja la variable
con el tipo \lstinline|?|.  Por esta razón, el chequeo estático
\lstinline|unknown :: Real| no falla, porque \lstinline|?| es
consistente con \lstinline|Real|.  Sin embargo, cuando se realiza el
chequeo dinámico de la ascripción en tiempo de ejecución, la
ascripción falla debido a la incompatibilidad del valor con el tipo
Real.

\begin{gplcblock}[caption={Ascripción de tipo \texttt{Real}},label={lst:ascr1}]
1 :: Real

>>>

{
	 [Real] 1. :: Real : SymProb1 
}  

\end{gplcblock}


\begin{gplcblock}[caption={Ascripción fallida},label={lst:ascr2}]
true :: Real

>>>

Static Error [line 1, cols 0-12]: In simple ascription, type Bool of term inconsistent with ascription type Real
\end{gplcblock}

\begin{gplcblock}[caption={Ascripción rechazada dinámicamente},label={lst:ascr3}]
let unknown = (true :: ?) in
(unknown :: Real)

>>>

Run-Time Error [line 2, cols 1-8]: expected value Real; found Bool
\end{gplcblock}


En los Códigos \ref{lst:ascr4} y \ref{lst:ascr5} se demuestra
de otra forma el poder de la ascripción para
modificar el tipo de una expresión. En el Código
\ref{lst:ascr4}, el tipo del parámetro es \lstinline|?|. Dada la
consistencia del tipo \lstinline|Bool| con \lstinline|?|, \lstinline|(id false)|
se ejecuta sin problemas. Sin embargo, en el
Código \ref{lst:ascr5},
el tipo de \lstinline|id| se \emph{refina} ascribiendo el tipo
de función \lstinline|Real -> {?:?}|, con dominio \lstinline|Real|.
La aplicación \lstinline|(real_id false)| falla estáticamente
por la inconsistencia del argumento con \lstinline|Real|,
demostrando el refinamiento del tipo esperado
para el argumento.


\begin{gplcblock}[caption={Lambda con parámetro sin tipo},label={lst:ascr4}]
let id = lambda x : ? . x in
(id false) 

>>>

{
	 [Bool] false :: ? : SymProb31 
}  
\end{gplcblock}





\begin{gplcblock}[caption={Lambda con parámetro de tipo refinado a \texttt{Real}},label={lst:ascr5}]
let id = lambda x : ? . x in
let real_id = id :: Real -> {?:?} in
(real_id false)

>>>

Static Error [line 3, cols 0-15]: In function, expected argument of type Real, but found Bool
\end{gplcblock}




\subsection{Integridad semántica probabilística}

El intérprete tiene la capacidad para juzgar semánticamente:
{(1)} bien o mal formuladas las probabilidades graduales usadas
en la selección probabilística y en la construcción de tipos de
distribución, y {(2)} bien o mal construidas las distribuciones
de tipos usadas en la ascripción.

En el caso de las probabilidades graduales, la probabilidad no
especificada \lstinline|?| es siempre correctamente utilizada. La
probabilidad numérica, en cambio, debe ser un número real en el intervalo $[0,1]$. En el Código \ref{lst:perr1}, la selección es rechazada estáticamente
debido al valor de la probabilidad.

\begin{gplcblock}[caption={Probabilidad inválida},label={lst:perr1}]
1.1 $ true : false

>>>

Static Error [line 1, cols 0-3]: Expected probability in [0,1] but found 1.1
\end{gplcblock}


En el caso de las distribuciones, la regla con respecto a la
integridad de las probabilidades tiene 3 casos: {(1)} Si la
distribución solo contiene probabilidades numéricas, estas deben sumar
exactamente 1.0. {(2)} Si solo hay probabilidades desconocidas,
la distribución es correcta. {(3)} Si la distribución combina
probabilidades numéricas y no especificadas, entonces la suma de las
probabilidades numéricas debe ser un número en el intervalo $[0,1]$.
La distribución del Código \ref{lst:perr2} es rechazada
estáticamente porque sus probabilidades numéricas suman más de 1.0.
La distribución del Código \ref{lst:perr3} falla con el mismo mensaje, pero esta vez debido a que sus probabilidades
numéricas suman menos de 1.0, y no hay ninguna probabilidad \lstinline|?| que pueda complementar la probabilidad restante.

\begin{gplcblock}[caption={Distribución con exceso de probabilidad},label={lst:perr2}]
(? $ true : 23) :: {Bool:0.5, Real:0.6, Real : ?}

>>>
Static Error [line 1, cols 19-49]: Invalid distribution type. The sum of probabilities must be 1.0.
If the sum of numeric probabilities is in [0, 1), the distribution must contain at least one probability '?' to complete it.
\end{gplcblock}

\begin{gplcblock}[caption={Distribución sin probabilidades suficientes},label={lst:perr3}]
(? $ true : 23) :: {Bool:0.4, Real:0.4}
\end{gplcblock}



%%%%%%%%%%%%%%%%%%%%%%%%%%%%%%%%%%%%%%%%%%%%%%%%%%%%%%%%
%%%%%%%%%%%%%%%%%%%%%%%%%%%%%%%%%%%%%%%%%%%%%%%%%%%%%%%%
\section{Limitaciones}

Esta sección tiene como objetivo presentar, de manera crítica, las
principales limitaciones identificadas durante el desarrollo del
intérprete de GPLC. En este trabajo, el intérprete se concibió como un
prototipo de investigación: su foco principal es materializar la
semántica del lenguaje, permitir la ejecución de programas y servir
como plataforma para experimentar con la gradualidad tanto en tipos
como en probabilidades. En consecuencia, existen aspectos del sistema
que se privilegiaron (corrección semántica y fidelidad a la
formulación teórica) por sobre otros (optimización o escalabilidad a
programas de mayor tamaño).


En las siguientes subsecciones, se organizan las limitaciones en
categorías que reflejan el tipo de dificultad observada: {(i)}
imprecisiones semánticas, {(ii)} evaluación de probabilidades
simbólicas en el resultado, {(iii)} eficiencia de la resolución de
fórmulas, {(iv)} rediseño de tipado para operaciones binarias,
{(v)} ausencia de mecanismos de igualdad para testing
semántico, y {(vi)} costos computacionales asociados a
\lstinline|Let| y \lstinline|Dascr| en el AST destino.


\subsection{Imprecisiones semánticas}

Durante la implementación y la experimentación con la ejecución de
programas se identificaron resultados o comportamientos no
esperados. Entre las causas, se encuentran errores o imprecisiones en
la especificación formal del lenguaje. Algunos de estos errores han
sido corregidos en una versión enviada a revisión para una nueva
versión de revista del trabajo de \citet{gplc}. Sin embargo, hay
aspectos que aún no han sido corregidos, y que, por lo tanto, se ven
reflejados en la implementación, que sigue dicha formulación de manera
literal.

Un ejemplo ilustrativo es el siguiente programa del Código
\ref{lst:bug}, donde se observa que la distribución final contiene
valores que no coinciden con la intuición acerca de la
ejecución. En el resultado, dos valores son compatibles con la
interpretación intuitiva: \lstinline|[Bool]true :: ?| y
\lstinline|[Real]13 :: Real|.  Sin embargo, aparecen además:
{(1)} un error dinámico \lstinline|ERROR(?)|, que indica que una
ascripción con tipo desconocido no pudo validarse en tiempo de
ejecución, y {(2)} un valor \lstinline|[Real]13 :: ?| que
indica que, en algún punto del pipeline, el valor $13$ fue ascrito con
tipo desconocido, a pesar de que en el Código \ref{lst:bug} no hay
ningún elemento sintáctico que haga la ascripción \lstinline|:: ?|
necesaria o esperable.

\begin{gplcblock}[caption={Distribución con valores no esperados},label={lst:bug}]
? $ (true :: ?) : 13

>>>

{
  [Bool] true :: ? : SymProb19 ,
  ERROR(?): SymProb20 ,
  [Real] 13. :: ? : SymProb21 ,
  [Real] 13. :: Real : SymProb22
}
\end{gplcblock}

La evidencia hallada durante el desarrollo indica que este tipo de
fenómeno se sitúa en la evaluación de la ascripción de distribución en
el AST destino. Dicha expresión tiene la estructura
\lstinline|Dascr(e,m,nu)|, donde \lstinline|e| es la evidencia,
\lstinline|m| es la subexpresión, y \lstinline|nu| es el tipo
ascrito. Primero se evalúa \verb|m| a una distribución de valores
\verb|V|, y luego se asigna a cada valor \verb|V| una evidencia y un
tipo de \verb|nu| para chequear consistencia en tiempo de
ejecución. En el proceso de asignar a cada valor de \verb|V| una
evidencia y un tipo de \lstinline|nu|, hay asignaciones que no se
hacen correctamente.  Esto ocurre porque la regla de evaluación no está manejando bien
la tolerancia que tiene el tipo \lstinline|?| a asociarse con valores
con los cuales es consistente (\lstinline|?| es \emph{consistente} con \lstinline|13|), pero a los cuales no fue ascrito en el
programa. El error no se repite en el Código \ref{lst:good-ascr},
con anotaciones de tipos concretos.

\begin{gplcblock}[caption={Ascripción correctamente evaluada.},label={lst:good-ascr}]
? $ (true :: Bool) : (13 :: Real)

>>>
{
	 [Bool] true :: Bool : SymProb19 ,
	 [Real] 13. :: Real : SymProb20 
}  
\end{gplcblock}

\subsection{Evaluación de probabilidades simbólicas}\label{subsec:lim-prob-simbolicas}

Actualmente, el intérprete retorna distribuciones simbólicas de
valores en las cuales las probabilidades están representadas por
variables simbólicas. Sin embargo, la fórmula asociada a la
distribución, que restringe los valores reales simultáneos que estas
variables pueden tomar, no se muestra en la salida.  Si bien las
funcionalidades que realizan la impresión de fórmulas están
implementadas en el código, el volumen de información contenido en una
fórmula correspondiente a un programa no trivial, suele ser demasiado
grande para entregar una información interpretable por el usuario, y
por esta razón se omite en la salida.  En consecuencia, el resultado
observado por el usuario corresponde esencialmente a un listado de
valores asociados a pesos variables, pero no a las restricciones
que hacen que el sistema sea soluble y semánticamente válido.

Esta decisión de diseño resulta útil en el contexto experimental del
prototipo, pues permite observar con claridad los valores que
contiene el soporte de una distribución, y verificar que un programa
termina exitosamente. No obstante, incluso en este contexto, resulta
valioso contar con mecanismos que expresen el resultado de manera
numérica, instanciando las probabilidades simbólicas cuando el usuario
lo requiera.


\subsection{Resolución incremental de fórmulas}\label{subsec:lim-incremental}

En el intérprete, la satisfacibilidad de fórmulas se utiliza de forma
transversal: es un paso esencial para determinar consistencia,
precisión, el meet de tipos de distribución, y para validar
restricciones heredadas o acumuladas en tiempo de ejecución. En la
implementación actual, la función \lstinline|solve| que conecta el
intérprete con el solucionador Z3 se comporta como una \textit{caja
  negra} que recibe fórmulas y retorna únicamente un valor booleano
(true si existe solución, false si no). Este diseño favorece una
integración simple del solucionador dentro del prototipo, pero impide
reutilizar trabajo entre chequeos sucesivos: cada invocación construye
un contexto nuevo del solucionador, vuelve a traducir la fórmula completa a
expresiones nativas de Z3 y a resolverla, encareciendo el cómputo
cuando estas verificaciones se repiten frecuentemente durante la
evaluación, lo cual ocurre generalmente.

Esta limitación se vuelve especialmente relevante en operaciones entre
distribuciones (por ejemplo, durante un \lstinline|meet|), donde la
fórmula a resolver se construye por acumulación de restricciones pues
se incluyen las fórmulas de las distribuciones argumento y se agregan
nuevas variables y ecuaciones asociadas al coupling. En este
escenario, aunque las subfórmulas heredadas puedan asumirse
satisfacibles (un invariante garantizado por el flujo del intérprete),
una solución para una subfórmula no necesariamente se extiende a una
solución del sistema global, porque la operación introduce
restricciones adicionales.


\subsection{Tipado de operaciones binarias}\label{subsec:lim-binops}

En la formulación implementada, las operaciones binarias aritméticas y
booleanas exigen que sus operandos sean consistentes con
\lstinline|Real| o \lstinline|Bool|, respectivamente. Una vez
satisfecha esta condición, el tipo de la operación se interpreta como
una distribución de Dirac sobre el tipo correspondiente (intuitivamente,
una distribución cuyo soporte contiene un único tipo con probabilidad
1). Este diseño es coherente con la idea de que toda expresión denota
una distribución. Sin embargo, genera una tensión en el sistema:
existen construcciones en el lenguaje (como aplicación de abstracción,
ascripción simple, operaciones binarias, y condiciones de un
\lstinline|if|) que esperan un término de tipo simple (un valor), no
un término cuyo tipo inferido sea una distribución (aunque sea de Dirac).


Como consecuencia, expresiones como \lstinline|(1+1)+7|,
\lstinline|(f (1 + 1))| o \lstinline|(1 + 1) :: Real| evidencian una
ambigüedad conceptual: la subexpresión \lstinline|(1 + 1)| puede verse
como un valor simple \lstinline|Real| o como una distribución
\lstinline|{Real:1}|. En el prototipo, esta tensión no se resolvió
mediante un rediseño del sistema de tipos, ya que el foco principal
del trabajo fue implementar la gradualidad en tipos de distribución y
en probabilidades desconocidas.


\subsection{Equivalencia semántica y testing}\label{subsec:lim-igualdad}

Una limitación importante para la validación sistemática es la
ausencia de un operador o mecanismo interno para comprobar igualdad
entre dos expresiones (o entre dos distribuciones resultantes). Sin
una noción de igualdad, resulta difícil automatizar tests que comparen
un resultado esperado con el que produce una expresión. Una
funcionalidad de este tipo tendría especial importancia para verificar
equivalencia semántica entre expresiones del lenguaje fuente,
comparando las distribuciones de valores que generan. Por ejemplo,
sería deseable poder contrastar directamente los programas:
\begin{itemize}
\item \lstinline|let x = (0.5 $ 1 : true) in 0.5 $ 1 : x| 
\item \lstinline|0.75 $ 1 : true|.
\end{itemize}


En el estado actual, lo más cercano a una comparación consiste en
ascribir manualmente un tipo de distribución esperado y verificar que
el programa termina sin error (Código \ref{lst:pseudo-comparison}), lo
que es una condición necesaria pero no suficiente para equivalencia
semántica.

\begin{gplcblock}[caption={Verificación de tipo de distribución},label={lst:pseudo-comparison}]
(let x = (0.5 $ 1 : true) in  0.5 $ 1 : x ) :: {Real : 0.75, Bool : 0.25}

>>>

{
	 [Real] 1. :: Real : SymProb95 ,
	 [Real] 1. :: Real : SymProb96 ,
	 [Real] 1. :: Real : SymProb97 ,
	 [Bool] true :: Bool : SymProb98 
}  
\end{gplcblock}



\subsection{Ineficiencia de \lstinline|Dascr| y \lstinline|Let|}\label{subsec:lim-eficiencia}

La evaluación de \lstinline|Let| y
\lstinline|Dascr| es un proceso computacionalmente
costoso.  Esto se debe a dos razones
estructurales. Primero, ambos constructos aparecen
con alta frecuencia en el AST destino: La
elaboración reestructura varios términos fuente
en \verb|Let|, y inserta ascripciones de distribución.  Segundo, su evaluación ejecuta
operaciones entre distribuciones de tipo que, en
el peor caso, crecen combinatoriamente: si dos
distribuciones de tipo con tamaños $m$
y $n$, se combinan mediante una operación como
meet, la distribución resultante puede
alcanzar tamaño $m \times n$ antes de descartar
combinaciones de tipos simples inválidas.

La implementación de las operaciones entre
distribuciones de tipos se basa en la construcción
de un coupling simbólico (una matriz de variables
frescas, con dimensión $m\times n$) cuyas variables
se usan para construir y resolver el sistema de
restricciones que valida la operación (ver sección
\ref{par:symbolic}). Dicho sistema, en conjunción
con las fórmulas heredadas de las distribuciones
sobre las que se opera, pasa a ser la fórmula de la
distribución resultante.  Este enfoque es
directo y cuida la claridad y fidelidad a la
definición, pero es muy costoso en tiempo y
memoria cuando se repite de manera profunda y
recurrente en la recursión estructural de la
evaluación sobre el AST destino.

%%%%%%%%%%%%%%%%%%%%%%%%%%%%%%%%%%%%%%%%%%%%%%%%%%%%%%%%
%%%%%%%%%%%%%%%%%%%%%%%%%%%%%%%%%%%%%%%%%%%%%%%%%%%%%%%%


\section{Resumen}

En conjunto, los casos presentados entregan evidencia práctica de que
el intérprete se comporta de manera consistente frente a los
escenarios estudiados: acepta programas bien formados y produce
resultados con la estructura esperada (distribuciones, sus valores y
las ascripciones asociadas), y rechaza programas inválidos cuando se
violan restricciones del lenguaje, reportando diagnósticos precisos
que incluyen localización y explicación de la inconsistencia
detectada. Además, los últimos ejemplos muestran que las reglas
semánticas sobre probabilidades válidas y distribuciones bien
construidas son efectivamente forzadas por el sistema.

Si bien esta validación no pretende ser exhaustiva, sí es sistemática
en el sentido de cubrir cada construcción central del lenguaje y una
colección de escenarios representativos y de borde. Por lo tanto, el
capítulo funciona tanto como una validación funcional del intérprete
como una guía concreta de uso del lenguaje GPLC, conectando
directamente las decisiones de diseño y las reglas semánticas con el
comportamiento observable de la implementación.

Finalmente, la subsección de Limitaciones sistematiza dificultades y
aspectos pendientes identificados durante el desarrollo del prototipo
(por ejemplo, imprecisiones semánticas puntuales, carácter poco
informativo de un resultado exclusivamente simbólico y costos asociados
al uso del solucionador). Estas observaciones acotan el alcance de la
validación y motivan las líneas de trabajo futuro sugeridas en el
capítulo final.
